% !PRINTED 168
\weberpage{168}
%<<CONT
worin $\lambda_1,\ldots\lambda_n$ Constanten sind, \"uber die wir so verf\"ugen, dass der Co\"efficient von $y^r$ in $F_0$ nicht verschwindet, d. h. dass $F_0(1,\lambda_1,\ldots\lambda_n)$ von Null verschieden wird. Dann geht das System $\mathfrak S$, wie \"uberhaupt jede Form der Variablen (10), in ein System von Formen der Variablen $y,x_1,\ldots x_n$ \"uber, und umgekehrt kann jede Form, die von diesen Variablen abh\"angt, auch als Form von den linearen Verbindungen (10) dargestellt werden.

Irgend eine Form $F$ des Systemes $\mathfrak S$ wird nun nach Potenzen von $y$ geordnet und dann in Bezug auf $y$ die Division mit $F_0$ ausgef\"uhrt, wobei sich
\weq{11}{F = a_0F_0 + \Phi}
ergeben mag, so dass $a_0$ der Quotient und $\Phi$ der Rest der Division ist. $a_0$ und $\Phi$ sind \gesperrt{ganze Functionen} der Variablen $x, y$, weil der Co\"efficient der h\"ochsten Potenz von $y$ im Divisor $F_0$ constant ist. $\Phi$ \"ubersteigt in Bezug auf $y$ nicht den Grad $r-1$. Durchl\"auft nun $F$ das System $\mathfrak S$, so bildet die Gesammtheit der durch (11) definirten Functionen $\Phi$ ein System $\mathfrak S_{r-1}$, von dem wir die Darstellbarkeit durch eine Basis als schon erwiesen annehmen. Wir k\"onnen also setzen:
\weq{12}{\Phi = A_1\Phi_1 + \cdots + A_\mu\Phi_\mu,}
so dass $\Phi_1,\ldots\Phi_\mu$ dem Systeme $\mathfrak S_{r-1}$ angeh\"oren, d. h. so, dass sich in dem Systeme $\mathfrak S$ die Formen
\weq{13}{F_1 = a_1F_0+\Phi_1,\ \ldots,\ F_\mu = a_\mu F_0+\Phi_\mu}
bestimmen lassen. Setzen wir also
\weq{14}{A_0 = a_0 - a_1A_1 - \cdots - a_\mu A_\mu,}
so folgt aus (11), (12) und (13):
\weq{15}{F = A_0F_0+A_1F_1+\cdots+A_\mu F_\mu,}
wodurch das Theorem I. allgemein bewiesen ist.

\begin{center}
\S. 42.\\
\gesperrt{Endlichkeit des Invariantensystemes einer endlichen linearen Substitutionsgruppe.}
\end{center}

Der im vorigen Paragraphen gegebene Beweis des Satzes I. ist an sich keinerlei Ausnahmen unterworfen. Wir verlieren aber nichts Wesentliches an seiner Allgemeinheit, wenn wir ein- f\"ur allemal identisch verschwindende Formen ausschliessen. Wenn
%CONT>>
